\section{MOSAIC}

\subsection{Evidence-Preserving Terrain Memory}

Depth pixels are back-projected, transformed into the current base frame, and splatted into
a measurement map $Z_t=[H_t^{z},V_t^{z}]$, containing mean height and empirical variance.
Previous memory is warped by planar odometry, $\bar{\vect m}_{t-1}=W(\vect m_{t-1},
\vect{x}_{t-1:t})$. A gated convolutional update then computes
\begin{align}
 \vect g_t &= \sigma\!\left(f_g[Z_t,\bar{\vect m}_{t-1},A_t,C_t]\right),\\
 \vect m_t &= \vect g_t\odot f_z(Z_t,C_t)
            +(1-\vect g_t)\odot\bar{\vect m}_{t-1},
 \label{eq:memory}
\end{align}
where $A_t$ is per-cell observation age and $C_t$ is a sparse contact map. If $d_t=0$,
$Z_t$ is masked and the update propagates memory while increasing age. Unlike a standard
ConvGRU, the gate receives explicit evidence age, preventing unchanged hidden state from
being interpreted as fresh evidence.

When foot $j$ establishes a stable contact, forward kinematics yields
$\tilde h_j$. We update nearby cells using a Gaussian kernel $k_j(i)$ and contact variance
$\sigma_c^2$:
\begin{equation}
 \hat h'_{t,i}=\frac{w_{t,i}\hat h_{t,i}+
 k_j(i)\sigma_c^{-2}\tilde h_j}{w_{t,i}+k_j(i)\sigma_c^{-2}}.
 \label{eq:contact}
\end{equation}
The same event reduces local risk only after slip and impact checks pass.

\subsection{Calibrated Risk Head}

Two decoder heads predict height $\hat H_t$ and risk logit $L_t$. Training sequences contain
random rectangular occluders, dropped frames, motion blur, and depth-dependent noise. The
risk label is one if the event in Eq.~\eqref{eq:risk} occurs within the cell. The loss is
\begin{equation}
\begin{aligned}
 \mathcal{L}={}&\lambda_h\lVert M\odot(\hat H-H^\star)\rVert_1
 +\lambda_n\mathcal{L}_{\mathrm{normal}}\\
 &+\lambda_r\operatorname{BCE}(L,Y)
 +\lambda_t\lVert\hat H_t-W(\hat H_{t-1})\rVert_1,
\end{aligned}
 \label{eq:loss}
\end{equation}
where $M$ marks valid ground truth. After training, a scalar temperature $\tau$ is fitted on
a held-out validation set and $R=\sigma(L/\tau)$. We report expected calibration error
(ECE) with 15 equal-mass bins and the Brier score; lower is better for both.

\subsection{Risk-Constrained Foothold Control}

The MPC stage extends a conventional centroidal objective with map terms:
\begin{equation}
 J=J_{\mathrm{track}}+\lambda_eJ_{\mathrm{effort}}
 +\sum_{k=1}^{K}\left[
 \lambda_s\lVert\nabla\hat H(\vect p_k)\rVert_2^2
 +\lambda_r\,\phi(r(\vect p_k))\right],
 \label{eq:objective}
\end{equation}
where $\phi(r)=-\log(1-r+10^{-4})$. Candidate contacts violating
Eq.~\eqref{eq:chance} are removed before continuous trajectory refinement. If no feasible
candidate remains for three replans, the robot lowers its body, shortens step length, and
turns toward the lowest-risk visible sector. This deterministic fallback avoids asking the
learned map to hallucinate a route through wholly unobserved space.

\begin{figure}[t]
\centering
\begin{tikzpicture}
\begin{axis}[
  width=\columnwidth,height=4.25cm,
  ybar,bar width=5.5pt,
  ymin=0,ymax=30,
  ylabel={Interventions (\%)},
  symbolic x coords={Nominal,Random,Structured,Dust},
  xtick=data,
  xticklabel style={font=\scriptsize,rotate=18,anchor=east},
  yticklabel style={font=\scriptsize},
  ylabel style={font=\small},
  legend style={font=\scriptsize,draw=none,fill=none,at={(0.5,1.02)},anchor=south,
                legend columns=3},
  axis line style={rssgray},
  grid=major,grid style={rssgray!15},
  enlarge x limits=0.16
]
\addplot[fill=rssgray!55,draw=rssgray] coordinates
  {(Nominal,5.0) (Random,14.6) (Structured,21.7) (Dust,27.1)};
\addplot[fill=rssteal!65,draw=rssteal] coordinates
  {(Nominal,4.6) (Random,10.2) (Structured,13.8) (Dust,18.5)};
\addplot[fill=rssorange!78,draw=rssorange] coordinates
  {(Nominal,4.2) (Random,6.5) (Structured,7.9) (Dust,11.3)};
\legend{Reactive,Memory-only,\method{}}
\end{axis}
\end{tikzpicture}
\caption{Illustrative hardware intervention rate across sensing conditions. Structured
dropout removes every fourth 0.8-s camera interval; dust uses a scattering overlay plus
real missing-depth masks.}
\label{fig:interventions}
\end{figure}
